% Derleme: XeLaTeX (Overleaf: Menu -> Compiler -> XeLaTeX)
\documentclass[11pt,a4paper]{article}
\usepackage{fontspec}
\usepackage{polyglossia}
\setdefaultlanguage{turkish}
\usepackage[margin=2.2cm]{geometry}
\usepackage{amssymb,amsmath}
\usepackage{enumitem}
\usepackage{booktabs,tabularx,longtable}
\usepackage[table]{xcolor}
\usepackage[most]{tcolorbox}
\usepackage{fancyhdr}
\usepackage{titlesec}
\usepackage[hidelinks]{hyperref}

\definecolor{ana}{HTML}{0F6E56}
\definecolor{acik}{HTML}{E1F5EE}
\definecolor{uyari}{HTML}{993C1D}
\definecolor{uyariacik}{HTML}{FAECE7}

\titleformat{\section}{\Large\bfseries\color{ana}}{\thesection}{0.6em}{}
\titleformat{\subsection}{\large\bfseries}{\thesubsection}{0.5em}{}
\setlength{\parindent}{0pt}
\setlength{\parskip}{4pt}
\setlength{\emergencystretch}{3em}\sloppy

\pagestyle{fancy}\fancyhf{}
\lhead{\small Bitirme Projesi -- Yapılacaklar Listesi}
\rhead{\small v1.4 -- 11 Kasım 2026}
\cfoot{\small\thepage}

% ---- Görev listesi ----
\newlist{gorev}{itemize}{2}
\setlist[gorev]{label=$\square$, leftmargin=1.6em, itemsep=5pt, parsep=1pt}
\newcommand{\cikti}[1]{\par{\small\textit{Çıktı:} #1}}
\newcommand{\kabul}[1]{\par{\small\textit{Kabul ölçütü:} #1}}
\newcommand{\donanim}{\textcolor{uyari}{\small[DONANIM]}\ }
\newcommand{\artı}{\textcolor{ana}{\small[FARK]}\ }

\newtcolorbox{faz}[2]{colback=acik,colframe=ana,fonttitle=\bfseries,
  title={#1},boxrule=0.6pt,arc=2pt,before upper={\textbf{Kilometre taşı:} #2\par}}
\newtcolorbox{kontrol}{colback=acik,colframe=ana,fonttitle=\bfseries,
  title={Haftalık kontrol},boxrule=0.6pt,arc=2pt}
\newtcolorbox{dikkat}{colback=uyariacik,colframe=uyari,boxrule=0.6pt,arc=2pt}

\begin{document}

\begin{center}
{\LARGE\bfseries\color{ana} EV Batarya Kontrol Ünitesi -- HIL Test Platformu}\\[4pt]
{\large Yapılacaklar Listesi ve Adım Adım Uygulama Planı}\\[4pt]
{\small 19S12P LG 21700 Dijital İkiz \textbullet{} EKF \textbullet{} Hafif ML \textbullet{} CAN \textbullet{} Güvenli Güç Kısma}
\end{center}

\section*{Nasıl kullanılır}
\begin{itemize}[leftmargin=1.2em,itemsep=1pt]
  \item Her görevin bir \textit{çıktısı} ve ölçülebilir bir \textit{kabul ölçütü} var. Ölçüt sağlanmadan kutu işaretlenmez.
  \item \donanim etiketli görevler fiziksel parça gerektirir. Diğer her şey parçalar gelmeden bilgisayarda yapılabilir.
  \item \artı etiketli görevler projeyi sıradan bir bitirme projesinden ayıran kısımlardır. Zaman sıkışırsa önce bunlar ertelenir, temel akış (MVP) değil.
  \item Tarihler 6 haftalık takvime göredir. Teslim tarihi kesinleşince haftalar kaydırılır, sıra değişmez.
\end{itemize}

% =====================================================================
\section{Proje mimarisi: 4 hedef, 4 modül}

\begin{center}
\renewcommand{\arraystretch}{1.25}
\begin{tabularx}{\textwidth}{>{\bfseries}p{3.2cm}X X}
\toprule
Hedef & Yazılım modülü & Fiziksel karşılığı \\
\midrule
1. Dijital ikiz ve taklit devre & 2-RC + 2 düğüm termal model, yaşlanma (stres) modeli, sürüş çevrimi & Emülatör Arduino Uno (Robotistan seti): ikizi 10 ms'de koşturur, ölçümleri CAN'a basar \\
2. Çift yönlü CAN & DBC, sayaç + CRC bütünlük koruması, zaman aşımı izleme & STM32 (dahili FDCAN + TJA1050) ile Arduino (MCP2515 + TJA1050), 500 kbps, uçlarda 120 $\Omega$ \\
3. Sinyal işleme + ML & Ön filtre, EKF (SOC), termal gözlemci ($T_\text{core}$), özellik çıkarımı, karar ağacı $\rightarrow \Psi_\text{stress}$ & ECU STM32 üzerinde, $<1$ ms döngü \\
4. Güvenli karar mekanizması & 4 durumlu makine (NORMAL, DERATE, LIMP-HOME, SAFE-STOP), histerezis, $I_\text{allow}$ & PWM $\rightarrow$ MOSFET $\rightarrow$ 12 V lamba \\
\bottomrule
\end{tabularx}
\end{center}

\textbf{Yazılım araçları (hepsi elde mevcut veya ücretsiz):} MATLAB/Simulink R2025b (Signal Processing ve DSP System Toolbox dahil), Python (numpy, scikit-learn), STM32CubeIDE, Arduino IDE, Git/GitHub, Overleaf.

\textbf{Donanım kararı (bütçe):} Tek STM32 (ECU) ve elde bulunan Arduino Uno (batarya emülatörü). Ağda iki farklı üreticinin CAN denetleyicisi (STM32 FDCAN/bxCAN ve Microchip MCP2515) konuşur; bu, protokolün standarda uygunluğunu gösterir ve gerçek araçlardaki çok tedarikçili ECU yapısına benzer.

\begin{dikkat}
\textbf{Lisans gerçeği ve çözümü.} Stateflow, Embedded Coder ve Statistics \& ML Toolbox yok. Bu yüzden:
(i) durum makinesi MATLAB Function bloğunda \texttt{switch-case} ile yazılır;
(ii) STM32 kodu STM32CubeIDE'de elle yazılır ve C kodu MATLAB çıktısıyla \emph{birebir} karşılaştırılır (SIL testi);
(iii) ML modeli Python'da eğitilir, C'ye aktarılır;
(iv) canlı izleme, Nucleo'nun ST-LINK USB'si üzerinden UART ile yapılır (MATLAB \texttt{serialport} veya Python).
Simulink'in harici modu Embedded Coder gerektirdiği için kullanılmıyor.
\end{dikkat}

% =====================================================================
\section{Faz 0 -- Altyapı ve sipariş (8--9 Kasım)}
\begin{faz}{Amaç: çalışma ortamını kurmak ve parçaları hemen sipariş etmek}{Simulink paketi çalışıyor, sipariş verildi}
Parçalara 4. haftaya kadar ihtiyaç yok, ama siparişi şimdi vermek gecikme riskini sıfırlar.
\end{faz}
\begin{gorev}
  \item GitHub'da özel bir repo aç. Klasör yapısı:
  \texttt{twin/ simulink/ ml/ firmware/common firmware/emulator firmware/ecu tools/ thesis/}
  \cikti{Repo bağlantısı; mevcut Python ve Simulink dosyaları yüklü}
  \kabul{Her çalışma günü en az bir commit atılıyor}
  \item[$\boxtimes$] Simulink paketini çalıştır: \texttt{build\_bms\_model}, ardından \texttt{run\_and\_compare}.
  \cikti{\texttt{bms\_twin\_ekf.slx} ve iki grafik}
  \kabul{Simulink ile Python arasındaki $V_\text{pack}$ farkı $<1$ mV; EKF RMSE $<\%1.5$}
  \item[$\boxtimes$] MATLAB'a MinGW-w64 eklentisini kur (Add-Ons, ücretsiz). PC'de C derlemek ve SIL testi için gerekli.
  \kabul{\texttt{mex -setup C} komutu MinGW'yi gösteriyor}
  \item[$\boxtimes$] \donanim Bölüm~\ref{sec:bom}'deki listeyi sipariş et.
  \kabul{Sipariş numaraları bu belgeye not edildi}
  \item Hocayla 15 dakikalık bir toplantı yap: kapsamı (4 hedef), teslim tarihini ve tez formatını netleştir.
  \cikti{Teslim tarihi ve tez şablonu}
\end{gorev}

% =====================================================================
\section{Hafta 1 -- Dijital ikiz ve yaşlanma modeli (9--15 Kasım)}
\begin{faz}{Hedef 1: tam fonksiyonlu, senaryo koşturabilen dijital ikiz}{senaryo grafikleri}
\end{faz}
\begin{gorev}
  \item[$\boxtimes$] Hücre parametrelerini LG M50T veri sayfası ve literatür değerleriyle gözden geçir; her parametrenin kaynağını tabloya yaz. (Hücre etiketi doğrulandı: INR21700 M50T.)
  \cikti{\texttt{bms\_param\_struct.m} içinde kaynak yorumları; tez için parametre tablosu}
  \item[$\boxtimes$] Senaryo kütüphanesi kur: (a) normal 25\,°C; (b) sıcak gün 40\,°C; (c) soğutma arızası ($R_u\times3$); (d) agresif sürüş; (e) derin deşarj; (f) uzun şarj.
  \cikti{\texttt{scenarios.m}: tek komutla 6 senaryo ve otomatik grafik}
  \kabul{(c) senaryosunda $T_\text{core}$ 45\,°C üstüne çıkıyor (derating testi için gerekli)}
  \item[$\boxtimes$] Sürüş çevrimini pedal girişine bağla: pedal (\%0--100) $\rightarrow$ akım talebi $\rightarrow$ ECU'dan gelen $I_\text{allow}$ ile sınırlama. Kapalı çevrimin ilk hali.
  \kabul{$I_\text{allow}$ düşünce paket akımı da düşüyor}
  \item[$\boxtimes$] \artı Zayıf hücre modeli: 19 hücreden birine \%10 düşük kapasite ve \%20 yüksek $R_0$ ver; ortalama ve zayıf hücre gerilimlerini ayrı çıkar.
  \kabul{Deşarj sonunda zayıf hücre $V_\text{min}$'e ortalamadan önce ulaşıyor}
  \item[$\boxtimes$] Yaşlanma (stres) modeli (Millner 2010 ve Pelletier 2017 temelli):
  takvim yaşlanması $\propto \exp(-E_a/RT)\cdot f(\text{SOC})$, çevrim yaşlanması $\propto f(\text{DoD}, T, C\text{-rate})$.
  Çevrim sayımı için Signal Processing Toolbox'taki \texttt{rainflow} fonksiyonunu kullan.
  \cikti{\texttt{aging\_model.m}: senaryo başına kapasite kaybı [\%/1000 km]}
  \kabul{Sıcak ve agresif senaryolar normalden belirgin şekilde daha fazla yaşlanma veriyor; değerler literatür mertebesinde}
\end{gorev}

% =====================================================================
\section{Hafta 2 -- Sinyal işleme ve kestirim (16--22 Kasım)}
\begin{faz}{Hedef 3, birinci yarı: gürültülü ölçümden güvenilir durum bilgisi}{karşılaştırma tablosu}
\end{faz}
\begin{gorev}
  \item[$\boxtimes$] Sensör hata kütüphanesi: gürültü, kayma, kuantalama (mevcut), bunlara ek olarak ani sıçrama (spike), donmuş sensör ve kopuk NTC.
  \cikti{\texttt{fault\_injector.m}, \texttt{sensor\_diag.m}, EKF inovasyon kapısı; 6 kabul testi geçti}
  \item[$\boxtimes$] Ön filtre: Signal Processing Toolbox ile akım ve sıcaklık için düşük geçiren filtre tasarla (\texttt{designfilt}); gecikme ve gürültü azaltma ödünleşimini raporla.
  \cikti{\texttt{prefilter\_study.py}, \texttt{prefilter\_design.m}: EKF girdisi filtrelenmez; eşikler için 2. derece Butterworth 0,5 Hz}
  \kabul{(Eski ölçüt ``gecikme $<50$ ms'' fiziksel olarak anlamsızdı, revize edildi.) Isıl eşiklerde gecikme $<1$ s ve titreşim $\ge\%90$ azalma (sağlandı: 0,45 s, \%97)}
  \item[$\boxtimes$] Kestirim yöntemlerini karşılaştır: Coulomb sayımı, OCV tablosu, EKF (sabit $R_0$), EKF ($R_0(T)$), EKF + akım kayması durumu.
  \cikti{Tez için tablo: RMSE, maksimum hata, yakınsama süresi, işlem yükü}
  \kabul{En iyi EKF tüm senaryolarda RMSE $<\%1.5$ (sağlandı: \%0,23--0,55; sabit $R_0$ 40\,°C'de \%2,3 ile aşıyor)}
  \item[$\boxtimes$] Termal gözlemci: 2 düğümlü modeli $T_\text{surf}$ ölçümüyle düzelten Kalman filtresi kur ve $T_\text{core}$'u kestir.
  \kabul{Soğutma arızası senaryosunda $|T_\text{core}-\hat T_\text{core}|<1$\,°C (sağlandı: 0,28\,°C; soğutma arızası 471 s'de algılandı)}
\item[$\boxtimes$] Çevrimiçi iç direnç kestirimi (ML'in sağlık göstergesi): $|\Delta I|\ge20$ A basamaklarında $\Delta V/\Delta I$, sıcaklık ve SOC'ye normalize, en küçük kareler (Das \& Kumar 2023 iHI fikri).
  \cikti{\texttt{rint\_estimator.m}; 3 sürüşte hata $<\%1$}
\end{gorev}

% =====================================================================
\section{Hafta 3 -- Karar mekanizması (23--29 Kasım)}
\begin{faz}{Hedef 3 ikinci yarı + Hedef 4}{kapalı çevrimde güç kısma simülasyonu}
\end{faz}
\subsection*{3a. Hafif makine öğrenmesi (2. haftaya öne çekildi, kapsam değişti)}
\textbf{Karar (11 Kasım):} Vekil (surrogate) stres modeli yerine \emph{gerçek veriyle SOH tahmini}: Kirkaldy vd. (2024) LG M50T verisi (40 hücre, 15 koşul). Etiketler gerçek ölçüm, kendi modelimizin çıktısı değil.
\begin{gorev}
  \item[$\boxtimes$] Veri hazırlama ve aykırı değer temizliği (\texttt{ml/01}); ilk modeller ve aşırı öğrenme (\texttt{ml/02}).
  \item[$\boxtimes$] Çapraz doğrulama \emph{koşul bazında} (leave-one-condition-out); veri sızıntısı karşılaştırması; fizik modeli, doğrusal, ağaç, gradient boosting, hibrit (\texttt{ml/03}).
  \kabul{Görülmemiş koşulda SOH RMSE $<\%3.5$ (sağlandı: \%2,99; hedef ML problemi değişince revize edildi)}
  \item[$\boxtimes$] Fizik kısıtlı karar ağacı (derinlik 3, monoton), C ve MATLAB'a otomatik aktarım (\texttt{ml/04}).
  \kabul{C ile Python aynı girdide $10^{-5}$ içinde eşit (sağlandı: $<4\cdot10^{-7}$)}
  \item[$\boxtimes$] Uçtan uca demo: çevrimiçi direnç $\rightarrow$ ML SOH $\rightarrow$ EKF, yaşlanmış pakette (\texttt{ml/05}).
  \item Simulink karşılığını çalıştır: \texttt{test\_ml\_soh}.
  \kabul{ML ile güncellenen BMS'in SOC RMSE'si yeni-hücre BMS'in en fazla yarısı}
\end{gorev}

\subsection*{3b. Güvenli karar mekanizması}
\begin{gorev}
  \item Durum makinesi (MATLAB Function, \texttt{switch-case}):
  \begin{itemize}[leftmargin=1.2em,itemsep=0pt]
    \item \textbf{NORMAL}: $I_\text{allow}=I_\text{nom}$
    \item \textbf{DERATE}: $I_\text{allow}=I_\text{nom}(1-\Psi_\text{stress})$, rampalı geçiş; $\Psi_\text{stress}$ çok kriterli: $\hat T_\text{core}$, $V_\text{cell,min}$, SOC, $\hat R_u$ alarmı, SOH, sensör bayrakları ve veri sayfasının sıcaklığa bağlı akım sınırları
    \item \textbf{LIMP-HOME}: $I_\text{allow}=I_\text{min}$ (lamba \%25, asla sönmez)
    \item \textbf{SAFE-STOP}: $I_\text{allow}=0$; $T>60$\,°C, $V_\text{cell}>4.25$ V, $V_\text{cell}<2.7$ V veya CAN zaman aşımı. Kilitlenir, sadece reset ile çıkılır.
  \end{itemize}
  \cikti{\texttt{decision\_logic.m} ve tez için durum diyagramı}
  \item Her geçişe histerezis ve süre filtresi ekle (örn. 40\,°C'de gir, 37\,°C'de 5 s bekleyerek çık).
  \kabul{Sınır değerin etrafında salınan girdide durum titreşmiyor}
  \item \artı Test matrisi: her durum geçişi için bir test senaryosu, beklenen sonuç ve gerçekleşen sonuç (basit FMEA tablosu).
  \kabul{Tüm geçişler en az bir testle kapsanıyor; tüm testler geçiyor}
  \item Kapalı çevrim simülasyonu: ikiz $\rightarrow$ sensör $\rightarrow$ EKF/ML $\rightarrow$ karar $\rightarrow$ akım sınırı $\rightarrow$ ikiz.
  \kabul{Soğutma arızasında $T_\text{core}$ derating sayesinde 55\,°C altında kalıyor; derating olmadan aşıyor (karşılaştırmalı grafik)}
\end{gorev}

% =====================================================================
\section{Hafta 4 -- Gömülü C, SIL ve CAN (30 Kasım -- 6 Aralık)}
\begin{faz}{Hedef 2 + algoritmaların karta taşınması}{iki kart CAN'da konuşuyor}
\end{faz}
\begin{gorev}
  \item \texttt{firmware/common/}: \texttt{ekf.c}, \texttt{thermal\_obs.c}, \texttt{stress\_tree.c}, \texttt{decision.c}, \texttt{bms\_can.c}. \texttt{float32}, dinamik bellek yok, platformdan bağımsız.
  \item \textbf{SIL testi (kart gerekmez):} ortak C kodunu PC'de derle, MATLAB'ın ürettiği girdi CSV'sini işlet, çıktıyı MATLAB ile karşılaştır.
  \kabul{SOC farkı $<0.05$ puan, durum geçişleri birebir aynı}
  \item CAN protokolü (DBC dosyası):
  0x101 Emülatör$\rightarrow$ECU ($V$, $I$, $T_\text{surf}$, $V_\text{zayıf}$),
  0x201 ECU$\rightarrow$Emülatör (SOC, $\hat T_\text{core}$, $\Psi$, $I_\text{allow}$, durum),
  0x301 komut/hata enjeksiyonu. Her mesajda 4 bitlik \emph{alive counter} ve CRC-8 (SAE J1850).
  \cikti{\texttt{bms.dbc} ve \texttt{bms\_can.c} kodlama/çözme fonksiyonları}
  \item \artı Zaman aşımı izleyici: 0x101 mesajı 50 ms gelmezse hata, 100 ms gelmezse SAFE-STOP.
  \item \donanim ECU: CubeIDE proje iskeleti (CAN, TIM PWM 1 kHz, UART, DWT sayacı).
  \item \donanim Emülatör: Arduino Uno + MCP2515 (CAN kütüphanesi, A0 pedal, buton kesmeleri, UART). \textbf{MCP2515 modülündeki kristali kontrol et (genelde 8 MHz) ve kütüphanede aynısını seç; yanlış seçilirse CAN hiç çalışmaz.}
  \item \donanim Önce her iki düğümü kendi geri döngü (loopback) modunda ayrı ayrı test et (STM32 CAN ve MCP2515'te var; ikinci düğüm olmadan ACK alınamaz), sonra iki düğümü hatta bağla; CAN-H ile CAN-L arası $\approx60$ $\Omega$ ölçülmeli.
  \kabul{10 dakikada 0 kayıp paket, hata sayaçları TEC = REC = 0}
  \item \donanim Lojik analizörü transceiver'ın RX pinine (3,3 V tarafı) bağla; PulseView'de CAN çözücüsünü 500 kbps'ye ayarla, 0x101 ve 0x201 çerçevelerini yakala ve tez için ekran görüntüsü al.
  \kabul{ID, veri baytları ve CRC çözülüyor; hata çerçevesi yok}
  \item Firmware içi CAN istatistikleri: her düğüm alive counter atlamalarını, CRC hatalarını, TEC/REC değerlerini sayar ve UART'tan PC'ye raporlar (10 dakikalık sıfır kayıp testinin kanıtı).
  \item \artı Kasıtlı hata üretimi emülatör firmware'inde: buton veya 0x301 komutuyla bozuk CRC, donmuş alive counter ve mesaj kesintisi üret; ECU'nun tepkisini doğrula.
\end{gorev}

% =====================================================================
\section{Hafta 5 -- HIL entegrasyonu ve canlı demo (7--13 Aralık)}
\begin{faz}{Tüm sistem masada, kapalı çevrimde}{demo videosu}
\end{faz}
\begin{gorev}
  \item \donanim Arduino'ya dijital ikizin C portunu yükle; 10 ms'lik periyotta çalıştır. Uno'da \texttt{double} = 32 bit olduğundan yük sayacı tamsayı (mA$\cdot$s) tutulur.
  \kabul{Kart çıktısı Simulink ile aynı çevrimde örtüşüyor (çapraz doğrulama 3: Python, Simulink ve C)}
  \item \donanim ECU döngü süresini DWT sayacıyla ölç.
  \kabul{$<1$ ms (hedef $<200\,\mu$s)}
  \item \donanim Far sürücüsü: STM32 PWM $\rightarrow$ IRLB8721 $\rightarrow$ 12 V LED şerit (veya setteki LED); gate'e 100 $\Omega$, gate-source arası 10k. Duty $= I_\text{allow}/I_\text{nom}$, alt sınır \%25.
  \item \donanim Hata enjeksiyonu butonları: soğutma arızası, sensör kopması, CAN kablosunu çekme.
  \kabul{Her biri doğru duruma geçiyor; CAN çekilince 100 ms içinde SAFE-STOP}
  \item Canlı izleme paneli: UART (ST-LINK sanal COM) $\rightarrow$ MATLAB \texttt{serialport} veya Python ile SOC, $T$, $\Psi$, durum ve $I_\text{allow}$ grafikleri.
  \item 2--3 dakikalık demo videosu çek (jüri günü donanım sorununa karşı B planı).
\end{gorev}

% =====================================================================
\section{Hafta 6 -- Tez ve savunma (14--20 Aralık)}
\begin{gorev}
  \item Tez bölümleri: Giriş ve literatür (5 makale) $\rightarrow$ Modelleme $\rightarrow$ Kestirim ve ML $\rightarrow$ Karar mekanizması $\rightarrow$ CAN ve HIL mimarisi $\rightarrow$ Sonuçlar (KPI tablosu) $\rightarrow$ Sonuç ve gelecek çalışmalar.
  \item Her KPI için bir grafik veya tablo; Python, Simulink ve C çapraz doğrulama tablosu.
  \item 10 dakikalık sunum (12--14 slayt) ve canlı demo provası; 30 saniyelik açılış konuşması.
  \item Olası jüri soruları ve cevapları (ML neden vekil model, emülatör neden sensör seviyesinde, güvenlik nasıl garanti ediliyor).
\end{gorev}

% =====================================================================
\section{Başarı ölçütleri (KPI)}
\begin{center}
\renewcommand{\arraystretch}{1.2}
\begin{tabularx}{\textwidth}{X l l}
\toprule
Metrik & Hedef & Doğrulama \\
\midrule
Simulink ve Python ikizi farkı & $<1$ mV & \texttt{run\_and\_compare} \\
EKF SOC RMSE & $<\%1.5$ & 6 senaryo, gürültü + kayma \\
EKF yakınsaması (\%50 $\rightarrow$ \%80) & $<60$ s & Yanlış başlangıç testi \\
$T_\text{core}$ kestirim hatası & $<1$\,°C & Soğutma arızası senaryosu \\
ML SOH doğruluğu & RMSE $<\%3.5$ & Leave-one-condition-out (gerçek veri) \\
C ile MATLAB farkı (SIL) & SOC $<0.05$ puan & PC'de derlenmiş C \\
ECU döngü süresi & $<1$ ms (hedef 200\,$\mu$s) & DWT sayacı \\
CAN güvenilirliği & 0 kayıp / 10 dk & Firmware sayaçları + TEC = REC = 0 \\
CAN zaman aşımı tepkisi & $<100$ ms $\rightarrow$ SAFE-STOP & Kablo çekme testi \\
Limp-home tabanı & Duty $\ge\%25$ & Isıl stres testi \\
Derating etkisi & $T_\text{core}<55$\,°C & Kapalı çevrim, derating var/yok \\
\bottomrule
\end{tabularx}
\end{center}

% =====================================================================
\section{Malzeme listesi ve ihtiyaç zamanı}\label{sec:bom}
\begin{center}
\renewcommand{\arraystretch}{1.2}
\begin{tabularx}{\textwidth}{X c l}
\toprule
Parça & Adet & Gerektiği hafta \\
\midrule
\textbf{Satın alınacak (zorunlu)} -- durum & & \\
NUCLEO-G431RB (Robotistan, \textbf{sipariş verildi}) & 1 & H4 \\
TJA1050 CAN modülü, STM32 tarafı (\textbf{sipariş verildi}, 11 Kasım). VCC $\rightarrow$ 5V pini; RX hattına 2,2k/3,3k gerilim bölücü & 1 & H4 \\
MCP2515 + TJA1050 CAN modülü, Arduino tarafı (Robotistan, \textbf{sipariş verildi}) & 1 & H4 \\
Far anahtarı: setteki NPN transistör (2N2222) veya mantık seviyeli MOSFET (IRLB8721) & 1 & H5 \\
Veri taşıyan Micro-USB kablosu (Nucleo için) & 1 & H4 \\
\midrule
\textbf{Robotistan Arduino setinden} & & \\
Arduino Uno R3 (CH340; Windows'ta CH340 sürücüsü gerekir), adaptör, breadboard, jumper kablolar & 1 & H4 \\
Potansiyometre (pedal), butonlar, dirençler, LED'ler; varsa 16x2 LCD (emülatör göstergesi) & -- & H4 \\
\midrule
\textbf{İsteğe bağlı} & & \\
Kısa 12 V LED şerit (yerine setteki LED'lerle far grubu yapılabilir) & 1 & H5 \\

USB lojik analizör, 24 MHz 8 kanal (PulseView ile CAN/UART/PWM çözme) & 1 & H4 \\
10k NTC ve 10 W direnç (gerçek ısıtma demosu) & 1 & H5 \\
\bottomrule
\end{tabularx}
\end{center}
Not: LG 21700 hücreler masada sadece görsel unsur olarak durur; hiçbir test gerçek hücre üzerinde yapılmaz.

% =====================================================================
\section{Riskler ve B planları}
\begin{center}
\renewcommand{\arraystretch}{1.2}
\begin{tabularx}{\textwidth}{X X}
\toprule
Risk & Önlem / B planı \\
\midrule
Kartlar geç gelir & H1--H3 ve SIL tamamen kartsız; HIL bir hafta kayabilir \\
CAN çalışmaz (bit zamanlaması, sonlandırma) & Önce loopback modu; lojik analizörle bit seviyesinde kontrol; 120 $\Omega$ ölçümü (hat üzerinde $\approx60$ $\Omega$) \\
ML doğruluğu düşük & Ağaç derinliğini artır veya özellik ekle; en kötü durumda kural tabanlı $\Psi$ ile devam (ML bonus olur) \\
Zaman yetmez & MVP sırası: ikiz $\rightarrow$ EKF $\rightarrow$ karar $\rightarrow$ CAN $\rightarrow$ lamba. \artı görevler en son \\
Demo günü donanım arızası & Hafta 5 videosu; emülatör yazılımı PC'de de koşabilir \\
\bottomrule
\end{tabularx}
\end{center}

% =====================================================================
\section{Literatür ve modül eşlemesi}
\begin{center}
\renewcommand{\arraystretch}{1.2}
\begin{tabularx}{\textwidth}{l X}
\toprule
Makale & Projede kullanıldığı yer \\
\midrule
Millner (2010) & Stres faktörlü yaşlanma modeli, rainflow çevrim sayımı (H1) \\
Pelletier vd. (2017) & Yaşlanma mekanizmaları, takvim ve çevrim ayrımı, paket ölçekleme (H1) \\
Ananthraj (2021) & BMS fonksiyonları, eşdeğer devre modeli, koruma eşikleri (H1, H3) \\
Das \& Kumar (2023) & Dolaylı sağlık göstergeleri (iHI), hafif ML yaklaşımı (H2, H3) \\
Xiong vd. (2019) & Gerçek BMS verisiyle sağlık prognozu, gömülü kısıtlar (H3, H4) \\
\bottomrule
\end{tabularx}
\end{center}

\vspace{6pt}
\begin{kontrol}
\begin{tabular}{@{}ll@{}}
$\boxtimes$ Faz 0: Simulink paketi çalıştı, parçalar sipariş edildi &
$\boxtimes$ H1: 6 senaryo ve yaşlanma modeli \\
$\boxtimes$ H2: Kestirim karşılaştırma tablosu &
$\square$ H3: Kapalı çevrimde derating \\
$\square$ H4: SIL geçti, CAN 0 kayıp &
$\square$ H5: Canlı demo ve video \\
$\square$ H6: Tez teslim, sunum provası & \\
\end{tabular}
\end{kontrol}

\end{document}
